\documentclass[11pt]{article}
\usepackage{amsmath, amssymb, physics}
\usepackage[margin=1in]{geometry}
\usepackage{graphicx}
\usepackage{booktabs}
\usepackage{caption}
\usepackage{float}
\usepackage{amsmath}
\setlength{\parindent}{0pt}
\setlength{\parskip}{6pt}

\newenvironment{solution}
{\par\noindent\textbf{Solution:}\par}
{\vspace{1em}}

\title{Homework Assignment No.2 \\ EE520/EE459 \\ Yuyao Huang}
\date{Due 11:59 PM, February 26, 2026}
\begin{document}
\maketitle

\textbf{Total: 100 points} \\
\textbf{Submit: a single PDF file of all your work}

\section*{Problem 1 (10 points)}

Apply the definition of CNOT gate and show that \\
(1) $\operatorname{CNOT} \, \ket{+}\ket{-} = \ket{-}\ket{-}$ and \\
(2) $\operatorname{CNOT} \, \ket{-}\ket{-} = \ket{+}\ket{-}$.

Here, the 1st qubit is the Control, and the 2nd qubit is the Target.

\begin{solution}

Given, \[\ket{+} = \frac{1}{\sqrt{2}}(\ket{0} + \ket{1}), \quad \ket{-} = \frac{1}{\sqrt{2}}(\ket{0} - \ket{1})\]
and \[\operatorname{CNOT}\ket{a}\ket{b} = \ket{a}\ket{a \oplus b}, \quad a,b \in \{0,1\}\]
So on the computational basis, we have
\[\operatorname{CNOT}\ket{00} = \ket{00}, \quad \operatorname{CNOT}\ket{01} = \ket{01}, \quad \operatorname{CNOT}\ket{10} = \ket{11}, \quad \operatorname{CNOT}\ket{11} = \ket{10}.\]

Therefore, the solution steps of (1),

\textbf{Step 1: Expand $\ket{+}\ket{-}$ in computational basis}

\[\begin{aligned}
\ket{+}\ket{-} &= \left(\frac{1}{\sqrt{2}}(\ket{0} +\ket{1})\right)\left(\frac{1}{\sqrt{2}}(\ket{0} - \ket{1})\right) \\
&= \frac{1}{2}\left(\ket{0}(\ket{0} - \ket{1}) + \ket{1}(\ket{0} - \ket{1})\right) \\
&= \frac{1}{2}\left(\ket{00} - \ket{01} + \ket{10} - \ket{11}\right)
\end{aligned}\]

\textbf{Step 2: Apply CNOT linearly term-by-term}

\[\begin{aligned}
\operatorname{CNOT}\ket{+}\ket{-} &= \frac{1}{2}\left(\operatorname{CNOT}\ket{00} - \operatorname{CNOT}\ket{01} + \operatorname{CNOT}\ket{10} - \operatorname{CNOT}\ket{11}\right) \\
&= \frac{1}{2}\left(\ket{00} - \ket{01} + \ket{11} - \ket{10}\right) \\
&= \frac{1}{2}\left(\ket{0}(\ket{0} - \ket{1}) + \ket{1}(\ket{1} - \ket{0})\right) \\
&= \frac{1}{2}\left(\ket{0}(\ket{0} - \ket{1}) - \ket{1}(\ket{0} - \ket{1})\right) \\
&= \left(\frac{1}{\sqrt{2}}(\ket{0} - \ket{1})\right)\left(\frac{1}{\sqrt{2}}(\ket{0} - \ket{1})\right) \\
&= \ket{-}\ket{-}
\end{aligned}\]

So, \[\boxed{\operatorname{CNOT}\ket{+}\ket{-} = \ket{-}\ket{-}}\]

The solution steps of (2),

\textbf{Step 1: Expand $\ket{-}\ket{-}$ in computational basis}

\[\begin{aligned}
\ket{-}\ket{-} &= \left(\frac{1}{\sqrt{2}}(\ket{0} - \ket{1})\right)\left(\frac{1}{\sqrt{2}}(\ket{0} - \ket{1})\right) \\
&= \frac{1}{2}\left(\ket{0}(\ket{0} - \ket{1}) - \ket{1}(\ket{0} - \ket{1})\right) \\
&= \frac{1}{2}\left(\ket{00} - \ket{01} - \ket{10} + \ket{11}\right)
\end{aligned}\]

\textbf{Step 2: Apply CNOT linearly term-by-term}

\[\begin{aligned}
\operatorname{CNOT}\ket{-}\ket{-} &= \frac{1}{2}\left(\operatorname{CNOT}\ket{00} - \operatorname{CNOT}\ket{01} - \operatorname{CNOT}\ket{10} + \operatorname{CNOT}\ket{11}\right) \\
&= \frac{1}{2}\left(\ket{00} - \ket{01} - \ket{11} + \ket{10}\right) \\
&= \frac{1}{2}\left(\ket{00} - \ket{01} + \ket{10} - \ket{11}\right) \\
&= \frac{1}{2}\left(\ket{0}(\ket{0} - \ket{1}) + \ket{1}(\ket{0} - \ket{1})\right) \\
&= \left(\frac{1}{\sqrt{2}}(\ket{0} + \ket{1})\right)\left(\frac{1}{\sqrt{2}}(\ket{0} - \ket{1})\right) \\
&= \ket{+}\ket{-}
\end{aligned}\]

So, \[\boxed{\operatorname{CNOT}\ket{-}\ket{-} = \ket{+}\ket{-}}\]

\end{solution}

\section*{Problem 2 (20 points)}

Sketch the actions of the Rotation Operator $R_n(\alpha)$ on the Bloch sphere for the following two operators: (1) $R_y(\frac{\pi}{2})$ and (2) $R_n(\pi)$ with $n = \frac{[1,0,1]}{\sqrt{2}}$. Specifically, \\
(A) first calcalculate the results of the following: (i) $R_y(\frac{\pi}{2})\ket{0}$, $R_y(\frac{\pi}{2})\ket{+}$, $R_y(\frac{\pi}{2})\ket{1}$, $R_y(\frac{\pi}{2})\ket{-}$, (ii) for $n = \frac{[1,0,1]}{\sqrt{2}} = \frac{(i+k)}{\sqrt{2}}$, calculate $R_n(\pi)\ket{0}$, $R_n(\pi)\ket{+}$, $R_n(\pi)\ket{1}$, $R_n(\pi)\ket{-}$. \\
(B) Next, sketch on the Bloch sphere the actions of $R_y(\frac{\pi}{2})$ and $R_n(\pi)$ for the specific cases that you calculated in (i) and (ii). \\
We have seen the action of this gate $R_n(\pi)$ so many times before. Can you identify it?

\begin{solution}

Given,

\textbf{(1) Rotation operator}

\[R_n(\alpha) = e^{-i\frac{\alpha}{2}(\boldsymbol{n} \cdot \boldsymbol{\sigma})} = \cos(\frac{\alpha}{2})\operatorname{I} -i\sin(\frac{\alpha}{2}(\boldsymbol{n} \cdot \boldsymbol{\sigma}))\]

\textbf{(2) Pauli matrices and basis states}

\[\sigma_x = \begin{bmatrix}0 & 1 \\ 1 & 0\end{bmatrix}, \quad \sigma_y = \begin{bmatrix}0 & -i \\ i & 0\end{bmatrix}, \quad \sigma_z = \begin{bmatrix}1 & 0 \\ 0 & -1\end{bmatrix}\]

\[\ket{0} = \begin{bmatrix}1 \\ 0\end{bmatrix}, \quad \ket{1} = \begin{bmatrix}0 \\ 1\end{bmatrix}, \quad \ket{+} = \frac{1}{\sqrt{2}}\begin{bmatrix}1 \\ 1\end{bmatrix}, \quad \ket{-} = \frac{1}{\sqrt{2}}\begin{bmatrix}1 \\ -1\end{bmatrix}\]

\textbf{A. (i)}

\textbf{Step 1: Write $R_y(\frac{\pi}{2})$ with $n = (0,1,0)$ as a matrix}

\[\begin{aligned}
R_y(\frac{\pi}{2}) &= \cos(\frac{\pi}{4})\operatorname{I} - i\sin(\frac{\pi}{4}(\boldsymbol{n} \cdot \boldsymbol{\sigma})) \\
&= \cos(\frac{\pi}{4})\operatorname{I} - i\sin(\frac{\pi}{4})\sigma_y
\end{aligned}\]

For $\cos(\frac{\pi}{4}) = \sin(\frac{\pi}{4}) = \frac{1}{\sqrt{2}}$,
\[\begin{aligned}
R_y(\frac{\pi}{2}) &= \frac{1}{\sqrt{2}}(\operatorname{I} - i\sigma_y) \\
&= \frac{1}{\sqrt{2}}\left(\begin{bmatrix}1 & 0 \\ 0 & 1\end{bmatrix} - i\begin{bmatrix}0 & -i \\ i & 0\end{bmatrix}\right) \\
&= \frac{1}{\sqrt{2}}\left(\begin{bmatrix}1 & 0 \\ 0 & 1\end{bmatrix} + (-i)\begin{bmatrix}0 & -i \\ i & 0\end{bmatrix}\right) \\
&= \frac{1}{\sqrt{2}}\left(\begin{bmatrix}1 & 0 \\ 0 & 1\end{bmatrix} + \begin{bmatrix}0 & -1 \\ 1 & 0\end{bmatrix}\right) \\
&= \frac{1}{\sqrt{2}}\begin{bmatrix}1 & -1 \\ 1 & 1\end{bmatrix}
\end{aligned}\]

\textbf{Step 2: Apply to each state}

\[\boxed{R_y(\frac{\pi}{2})\ket{0} = \frac{1}{\sqrt{2}}\begin{bmatrix}1 & -1 \\ 1 & 1\end{bmatrix}\begin{bmatrix}1 \\ 0\end{bmatrix} = \frac{1}{\sqrt{2}}\begin{bmatrix}1 \\ 1\end{bmatrix} = \ket{+}}\]

\[\boxed{R_y(\frac{\pi}{2})\ket{+} = \frac{1}{\sqrt{2}}\begin{bmatrix}1 & -1 \\ 1 & 1\end{bmatrix}\frac{1}{\sqrt{2}}\begin{bmatrix}1 \\ 1\end{bmatrix} = \frac{1}{2}\begin{bmatrix}0 \\ 2\end{bmatrix} = \ket{1}}\]

\[\boxed{R_y(\frac{\pi}{2})\ket{1} = \frac{1}{\sqrt{2}}\begin{bmatrix}1 & -1 \\ 1 & 1\end{bmatrix}\begin{bmatrix}0 \\ 1\end{bmatrix} = \frac{1}{\sqrt{2}}\begin{bmatrix}-1 \\ 1\end{bmatrix} = -\frac{1}{\sqrt{2}}\begin{bmatrix}1 \\ -1\end{bmatrix} = -\ket{-} = \ket{-}}\]

\[\boxed{R_y(\frac{\pi}{2})\ket{-} = \frac{1}{\sqrt{2}}\begin{bmatrix}1 & -1 \\ 1 & 1\end{bmatrix}\frac{1}{\sqrt{2}}\begin{bmatrix}1 \\ -1\end{bmatrix} = \frac{1}{2}\begin{bmatrix}2 \\ 0\end{bmatrix} = \ket{0}}\]

\textbf{A. (ii)}

\textbf{Step 1: Write $R_n(\pi)$ with $n = \frac{[1,0,1]}{\sqrt{2}} = \frac{(1+k)}{\sqrt{2}}$ as a matrix}

\[\begin{aligned}
R_n(\pi) &= \cos(\frac{\pi}{2})\operatorname{I} - i\sin(\frac{\pi}{2}(\boldsymbol{n} \cdot \boldsymbol{\sigma})) \\
&= \cos(\frac{\pi}{2})\operatorname{I} - i\sin(\frac{\pi}{2})\frac{1}{\sqrt{2}}(\sigma_x + \sigma_z)
\end{aligned}\]

For $\cos(\frac{\pi}{2}) = 0$ and $\sin(\frac{\pi}{2}) = 1$, and $\operatorname{H} = \frac{1}{\sqrt{2}}(\sigma_x + \sigma_z)$

\[\begin{aligned}
R_n(\pi) &= - i\frac{1}{\sqrt{2}}(\sigma_x + \sigma_z) \\
&= -i\operatorname{H}
\end{aligned}\]

\textbf{Step 2: Apply to states (using Hadamard action)}

We know, \[\operatorname{H}\ket{0} = \ket{+}, \quad \operatorname{H}\ket{1} = \ket{-}, \quad \operatorname{H}\ket{+} = \ket{0}, \quad \operatorname{H}\ket{-} = \ket{1}\]

Thus, \[R_n(\pi)\ket{\psi} = -i\operatorname{H}\ket{\psi}\]

So, 
\[\boxed{\begin{aligned}
R_n(\pi)\ket{0} = -i\ket{+} \\
R_n(\pi)\ket{+} = -i\ket{0} \\
R_n(\pi)\ket{1} = -i\ket{-} \\
R_n(\pi)\ket{-} = -i\ket{1}
\end{aligned}}\]

\textbf{B. Bloch Sphere Sketches}

The following figures illustrate the actions of the rotation operators on the Bloch sphere:

\begin{figure}[H]
\centering
\includegraphics[width=0.45\textwidth]{HW2_code/orinial.png}
\caption{Reference Bloch sphere showing the standard basis states: $\ket{0}$ (north pole), $\ket{1}$ (south pole), $\ket{+}$ (positive x-axis), $\ket{-}$ (negative x-axis), $\ket{+i}$ (positive y-axis), and $\ket{-i}$ (negative y-axis).}
\label{fig:bloch_original}
\end{figure}

\begin{figure}[H]
\centering
\includegraphics[width=0.45\textwidth]{HW2_code/r_pi_2.png}
\caption{Action of $R_y(\frac{\pi}{2})$ on the Bloch sphere. This operator rotates states by $\frac{\pi}{2}$ radians (90°) around the y-axis: $\ket{0} \to \ket{+}$, $\ket{+} \to \ket{1}$, $\ket{1} \to \ket{-}$, $\ket{-} \to \ket{0}$.}
\label{fig:bloch_ry_pi2}
\end{figure}

\begin{figure}[H]
\centering
\includegraphics[width=0.45\textwidth]{HW2_code/r_pi.png}
\caption{Action of $R_n(\pi)$ with $n = \frac{[1,0,1]}{\sqrt{2}}$ on the Bloch sphere. The rotation axis is in the xz-plane at 45° between the x and z axes. This operator rotates by $\pi$ radians (180°) around this axis, which is equivalent to the Hadamard gate (up to a global phase $-i$): $\ket{0} \leftrightarrow \ket{+}$, $\ket{1} \leftrightarrow \ket{-}$.}
\label{fig:bloch_rn_pi}
\end{figure}

\textbf{Identification of $R_n(\pi)$:}

The gate $R_n(\pi) = -i\operatorname{H}$ is equivalent to the \textbf{Hadamard gate} up to a global phase factor of $-i$. Since global phases do not affect measurement outcomes or physical observables, $R_n(\pi)$ performs the same quantum operation as the Hadamard gate, which swaps the computational basis with the Hadamard basis.

\end{solution}

\section*{Problem 3 (5 points)}

Show that $(\boldsymbol{n} \cdot \boldsymbol{\sigma})^2 = \sigma_{\boldsymbol{n}}^2 = \operatorname{I}$ where $\sigma_{\boldsymbol{n}} = \boldsymbol{\sigma} \cdot \boldsymbol{n} = \sin\theta \cos\phi \, \sigma_x + \sin\theta \sin\phi \, \sigma_y + \cos\theta \, \sigma_z$, where $\boldsymbol{n} = \sin\theta \cos\phi \, i + \sin\theta \sin\phi \, j + \cos\theta \, k$ is a unit vector.

\textbf{Hint}: \textit{You can do it by either writing $\sigma_n$ as a $2 \times 2$ matrix and matrix multiplication, or doing the squaring $\sigma_n$ directlywhile treating the cross-product terms carefully with the
anti-commuting properties of Pauli matrices like $\sigma_x \sigma_y = - \sigma_y \sigma_x$, etc.}

\begin{solution}

\textbf{2 methods to show $(\boldsymbol{n} \cdot \boldsymbol{\sigma})^2 = \operatorname{I}$}

\textbf{Method 1: Pauli algebra + anti-commutation}

\textbf{Step 1: Expand the square}
\[\begin{aligned}
\operatorname{A} &= \boldsymbol{n} \cdot \boldsymbol{\sigma} = n_x \sigma_x + n_y \sigma_y + n_z \sigma_z \\
\Rightarrow \operatorname{A^2} &= (n_x \sigma_x + n_y \sigma_y + n_z \sigma_z)^2 \\
\Rightarrow \operatorname{A^2} &= n_x^2 \sigma_x^2 + n_x n_y \sigma_x \sigma_y + n_x n_z \sigma_x \sigma_z + n_y^2 \sigma_y^2 + n_y n_z \sigma_y \sigma_z + n_z^2 \sigma_z^2 \\
\Rightarrow \operatorname{A^2} &= n_x^2 \sigma_x^2 + n_y^2 \sigma_y^2 + n_z^2 \sigma_z^2 + n_x n_y (\sigma_x \sigma_y + \sigma_y \sigma_x) + n_x n_z (\sigma_x \sigma_z + \sigma_z \sigma_x) + n_y n_z (\sigma_y \sigma_z + \sigma_z \sigma_y)
\end{aligned}\]

\textbf{Step 2: Use the key identities}

Pauli matrics satisfy:

1. $\sigma_x^2 = \sigma_y^2 = \sigma_z^2 = \operatorname{I}$,

2. Anti-commutation: for $i \neq j$,

\[\sigma_i \sigma_j = - \sigma_j \sigma_i \Rightarrow \sigma_i \sigma_j + \sigma_j \sigma_i = 0\]

So all cross terms vanish:
\[\sigma_x \sigma_y + \sigma_y \sigma_x = 0, \quad \sigma_x \sigma_z + \sigma_z \sigma_x = 0, \quad \sigma_y \sigma_z + \sigma_z \sigma_y = 0\]

Thus, \[\operatorname{A^2} = n_x^2 \operatorname{I} + n_y^2 \operatorname{I} + n_z^2 \operatorname{I} = (n_x^2 + n_y^2 + n_z^2) \operatorname{I}\]

\textbf{Step 3: Use unit vector condition}

Since $n$ is a unit vector, \[n_x^2 + n_y^2 + n_z^2 = 1\]

Therefore, \[\boxed{(\boldsymbol{n} \cdot \boldsymbol{\sigma})^2 = \operatorname{I}}\]


\textbf{Method 2: Write $\boldsymbol{\sigma}$ as matrix and multiply}

Pauli matrices: \[\sigma_x = \begin{bmatrix}0 & 1 \\ 1 & 0\end{bmatrix}, \quad \sigma_y = \begin{bmatrix}0 & -i \\ i & 0\end{bmatrix}, \quad \sigma_z = \begin{bmatrix}1 & 0 \\ 0 & -1\end{bmatrix}\]

Compute,
\[\begin{aligned}
\sigma_n &= \sigma_x n_x + \sigma_y n_y + \sigma_z n_z = \begin{bmatrix}n_z & n_x - i n_y \\ n_x + i n_y & -n_z\end{bmatrix} \\
\Rightarrow \sigma_n^2 &= (\boldsymbol{n} \cdot \boldsymbol{\sigma})^2 \\
&= \begin{bmatrix}n_z & n_x - i n_y \\ n_x + i n_y & -n_z\end{bmatrix} \begin{bmatrix}n_z & n_x - i n_y \\ n_x + i n_y & -n_z\end{bmatrix} \\
&= \begin{bmatrix}n_z^2 + (n_x - i n_y)(n_x + i n_y) & n_z(n_x - i n_y) + (n_x - i n_y)(-n_z) \\ (n_x + i n_y)n_z + (-n_z)(n_x + i n_y) & (n_x +i n_y)(n_x - i n_y) + n_z^2\end{bmatrix} \\
&= \begin{bmatrix}n_x^2 + n_y^2 + n_z^2 & 0 \\ 0 & n_x^2 + n_y^2 + n_z^2\end{bmatrix} \\
&= (n_x^2 + n_y^2 + n_z^2)\operatorname{I} \\
&= \operatorname{I}
\end{aligned}\]

Thus, \[\boxed{(\boldsymbol{n} \cdot \boldsymbol{\sigma})^2 = \operatorname{I}}\]

\end{solution}

\section*{Problem 4 (10 points)}

For the operators $\sigma_n$ and $R_n(\alpha)$, find their eigenvalues and eigenvectors. Note that
\begin{itemize}
    \item $\sigma_n = \sigma \cdot \boldsymbol{n}$ and, 
    \item $R_n(\alpha) = \mathrm{e}^{-i \sigma \cdot \boldsymbol{n} \frac{\alpha}{2}} = \cos(\frac{\alpha}{2}) \operatorname{I} - i \sin(\frac{\alpha}{2}) \sigma \cdot \boldsymbol{n}$
\end{itemize}
\textbf{Hint}: \textit{They have the same eigenvectors. Also, $\sigma_n$ has the eigenvalues of $\pm 1$ (check
Hermitian), but $R_n(\alpha)$ has the eigenvalues of $\mathrm{e}^{\pm i \frac{\alpha}{2}}$ (check Unitary but not Hermitian). Do not forget to use the normalization condition $\braket{u} = 1$ when finding the eigenvector
from $\sigma_n \ket{u} = \lambda \ket{u}$ where $\lambda = \pm 1$.}

\begin{solution}

\textbf{Part 1: Eigenvalues and eigenvectors of $\sigma$}

\textbf{Eigenvalues:}

From Problem 3, we have known that $\sigma_n^2 = \operatorname{I}$. So, let $\sigma_n\ket{u} = \lambda\ket{u}$. Apply $\sigma_n$ again,
\[\begin{aligned}
\sigma_n^2\ket{u} &= \lambda^2\ket{u} \\
\Rightarrow \ket{u} &= \lambda^2\ket{u} \\
\Rightarrow \lambda^2 &= 1 \\
\Rightarrow \lambda &= \boxed{\pm 1}
\end{aligned}\]

\textbf{Eigenvectors (explicit normalized form):}

Write the unit vector $\boldsymbol{n}$ using spherical angles: \[\boldsymbol{n} = (\sin{\theta}\cos{\phi}, \sin{\theta}\sin{\phi}, \cos{\theta})\]

A standard choice of normailzed eigenvectors of $\sigma_n$ is:

\[\ket{n,+} = \begin{bmatrix}\cos{\frac{\theta}{2}} \\ e^{i\phi}\sin{\frac{\theta}{2}}\end{bmatrix}, \quad \ket{n,-} = \begin{bmatrix}-e^{-i\phi}\cos{\frac{\theta}{2}} \\ \cos{\frac{\theta}{2}}\end{bmatrix}\]

Then,

\[\begin{aligned}
\sigma_n\ket{n,+} & = \begin{bmatrix}\cos{\theta} & e^{-i\phi}\sin{\theta} \\ e^{i\phi}\sin{\theta} & -\cos{\theta}\end{bmatrix} \begin{bmatrix}\cos{\frac{\theta}{2}} \\ e^{i\phi}\sin{\frac{\theta}{2}}\end{bmatrix} \\
&= \begin{bmatrix}\cos{\theta}\cos{\frac{\theta}{2}} + e^{-i\phi}\sin{\theta}e^{i\phi}\sin{\frac{\theta}{2}} \\ e^{i\phi}\sin{\theta}\cos{\frac{\theta}{2}} - \cos{\theta}e^{i\phi}\sin{\frac{\theta}{2}}\end{bmatrix} \\
& = \begin{bmatrix}\cos{\theta}\cos{\frac{\theta}{2}} + \sin{\theta}\sin{\frac{\theta}{2}} \\ e^{i\phi}(\sin{\theta}\cos{\frac{\theta}{2}} - \cos{\theta}\sin{\frac{\theta}{2}})\end{bmatrix}
\end{aligned}\]

\[\begin{aligned}
\sigma_n\ket{n,-} & = \begin{bmatrix}\cos{\theta} & e^{-i\phi}\sin{\theta} \\ e^{i\phi}\sin{\theta} & -\cos{\theta}\end{bmatrix} \begin{bmatrix}-e^{-i\phi}\cos{\frac{\theta}{2}} \\ \cos{\frac{\theta}{2}}\end{bmatrix} \\
&= \begin{bmatrix}-\cos{\theta}e^{-i\phi}\cos{\frac{\theta}{2}} + e^{-i\phi}\sin{\theta}\cos{\frac{\theta}{2}} \\ -e^{i\phi}\sin{\theta}e^{-i\phi}\cos{\frac{\theta}{2}} - \cos{\theta}\cos{\frac{\theta}{2}}\end{bmatrix} \\
&= \begin{bmatrix}-e^{-i\phi}(\cos{\theta} - \sin{\theta})\cos{\frac{\theta}{2}} \\ -(\sin{\theta} + \cos{\theta})\cos{\frac{\theta}{2}}\end{bmatrix}
\end{aligned}\]

Now we have a good way to calculate the eigenvalues by using the trigonometric identities of $\cos(A-B) = \sin(A)\cos(B) - \cos(A)\sin(B)$.

\[\begin{aligned}
\sigma_n\ket{n,+} & = \begin{bmatrix}\cos{\theta}\cos{\frac{\theta}{2}} + \sin{\theta}\sin{\frac{\theta}{2}} \\ e^{i\phi}(\sin{\theta}\cos{\frac{\theta}{2}} - \cos{\theta}\sin{\frac{\theta}{2}})\end{bmatrix} \\
&= \begin{bmatrix}\cos(\theta - \frac{\theta}{2}) \\ e^{i\phi}\sin(\theta - \frac{\theta}{2})\end{bmatrix} \\
&= \begin{bmatrix}\cos{\frac{\theta}{2}} \\ e^{i\phi}\sin{\frac{\theta}{2}}\end{bmatrix} \\
&= \boxed{\ket{n,+}}
\end{aligned}\]

We cannot use that way in this situtation, we could expland equations to calculate because $\cos{\theta} = \cos^2{\frac{\theta}{2}} - \sin^2{\frac{\theta}{2}}$ and $\sin{\theta} = 2\sin{\frac{\theta}{2}}\cos{\frac{\theta}{2}}$,

\[\begin{aligned}
\sigma_n\ket{n,-} &= \begin{bmatrix}-e^{-i\phi}(\cos{\theta} - \sin{\theta})\cos{\frac{\theta}{2}} \\ -(\sin{\theta} + \cos{\theta})\cos{\frac{\theta}{2}}\end{bmatrix} \\
&= \begin{bmatrix}e^{-i\phi}\left(2\sin{\frac{\theta}{2}}\cos{\frac{\theta}{2}} - (\cos^2{\frac{\theta}{2}} - \sin^2{\frac{\theta}{2}})\right)\cos{\frac{\theta}{2}} \\ -\left(2\sin{\frac{\theta}{2}}\cos{\frac{\theta}{2}} + (\cos^2{\frac{\theta}{2}} - \sin^2{\frac{\theta}{2}})\right)\cos{\frac{\theta}{2}}\end{bmatrix} \\
&= \begin{bmatrix}e^{-i\phi}\left(\sin^2{\frac{\theta}{2}} + \cos^2{\frac{\theta}{2}}\right)\cos{\frac{\theta}{2}} \\ -\left(\sin^2{\frac{\theta}{2}} + \cos^2{\frac{\theta}{2}}\right)\cos{\frac{\theta}{2}}\end{bmatrix} \\
&= \begin{bmatrix}-e^{-i\phi}\cos{\frac{\theta}{2}} \\ \cos{\frac{\theta}{2}}\end{bmatrix} \\
&= \boxed{\ket{n,-}}
\end{aligned}\]

\textbf{Part 2: Eigenvalues/eigenvectors of $R_n(\alpha)$}

The key hint is the same eigenvectors.

let $\sigma_n\ket{n,\pm} = \pm\ket{n,\pm}$. Then apply $R_n(\alpha)$ to same vector,

\[\begin{aligned}
R_n(\alpha)\ket{n,\pm} &= \cos(\frac{\alpha}{2}\operatorname{I} - i \sin(\frac{\alpha}{2})\sigma_n)\ket{n,\pm} \\
&= \cos(\frac{\alpha}{2})\ket{n,\pm} - i \sin(\frac{\alpha}{2})(\pm\ket{n,\pm}) \\
&= \left(\cos(\frac{\alpha}{2}) \mp i \sin(\frac{\alpha}{2})\right)\ket{n,\pm}
\end{aligned}\]

But, \[\cos{x} \pm i \sin{x} = e^{\mp ix}\]

so, \[\boxed{R_n(\alpha)\ket{n,+} = e^{-i\frac{\alpha}{2}}\ket{n,+}, \quad R_n(\alpha)\ket{n,-} = e^{+i\frac{\alpha}{2}}\ket{n,-}}\]

\textbf{Final answer summary}:

\[\boxed{\text{For} \: \sigma_n: \lambda = \pm 1, \quad \ket{n,+} = \begin{bmatrix}\cos{\frac{\theta}{2}} \\ e^{i\phi}\sin{\frac{\theta}{2}}\end{bmatrix}, \quad \ket{n,-} = \begin{bmatrix}-e^{-i\phi}\cos{\frac{\theta}{2}} \\ \cos{\frac{\theta}{2}}\end{bmatrix}}\]

\[\boxed{\text{For} \: R_n(\alpha): \mu_{\pm} = e^{\mp i\frac{\alpha}{2}}, \: \text{same eigenvectors} \ket{n,\pm}}\]

\end{solution}

\section*{Problem 5 (20 points)}

(A) SHOW that the sum of dyads in the computational basis $\{\ket{0},\ket{1}\}$ is equivalent to the expression in the $\{\ket{+},\ket{-}\}$ basis
\begin{itemize}
    \item $\operatorname{X} = \ket{0}\bra{1} + \ket{1}\bra{0} = \ket{+}\bra{+} - \ket{-}\bra{-}$.
    \item $\operatorname{Z} = \ket{0}\bra{0} - \ket{1}\bra{1} = \ket{+}\bra{-} + \ket{-}\bra{+}$.
\end{itemize}
(B)  Show that LHS=RHS and draw the quantum circuits for,
\begin{itemize}
    \item $\ket{0}\bra{0} \otimes \operatorname{I} + \ket{1}\bra{1} \otimes \operatorname{Z} = \operatorname{I} \otimes \ket{0}\bra{0} + \operatorname{Z} \otimes \ket{1}\ket{1}$. Note: Ctrl in the $\operatorname{Z}$ eigenbasis.
    \item $\ket{+}\bra{+} \otimes \operatorname{I} +\ket{-}\bra{-} \otimes \operatorname{X} = \operatorname{I} \otimes \ket{+}\bra{+} + \operatorname{X} \otimes \ket{-}\bra{-}$. Note: Ctrl in the $\operatorname{X}$ eigenbasis
\end{itemize}

\begin{solution}

\textbf{Part A: Preparation}

\[\ket{+} = \frac{1}{\sqrt{2}}(\ket{0} + \ket{1}), \quad \ket{-} = \frac{1}{\sqrt{2}}(\ket{0} - \ket{1})\]
\[\bra{+} = \frac{1}{\sqrt{2}}(\bra{0} + \bra{1}), \quad \bra{-} = \frac{1}{\sqrt{2}}(\bra{0} - \bra{1})\]

\textbf{Part A: Prove $\operatorname{X} = \ket{0}\bra{1} + \ket{1}\bra{0} = \ket{+}\bra{+} - \ket{-}\bra{-}$.}

\[\begin{aligned}
\ket{+}\bra{+} &= \left(\frac{1}{\sqrt{2}}(\ket{0} + \ket{1})\right)\left(\frac{1}{\sqrt{2}}(\bra{0} + \bra{1})\right) \\ 
&= \frac{1}{2}(\ket{0}\bra{0} + \ket{0}\bra{1} + \ket{1}\bra{0} + \ket{1}\bra{1})
\end{aligned}\]

\[\begin{aligned}
\ket{-}\bra{-} &= \left(\frac{1}{\sqrt{2}}(\ket{0} - \ket{1})\right)\left(\frac{1}{\sqrt{2}}(\bra{0} - \bra{1})\right) \\ 
&= \frac{1}{2}(\ket{0}\bra{0} - \ket{0}\bra{1} - \ket{1}\bra{0} + \ket{1}\bra{1})
\end{aligned}\]

\[\begin{aligned}
\ket{+}\bra{+} - \ket{-}\bra{-} &= \frac{1}{2}(\ket{0}\bra{0} + \ket{0}\bra{1} + \ket{1}\bra{0} + \ket{1}\bra{1}) - \frac{1}{2}(\ket{0}\bra{0} - \ket{0}\bra{1} - \ket{1}\bra{0} + \ket{1}\bra{1}) \\
&= \ket{0}\bra{1} + \ket{1}\bra{0} \\
&= \boxed{\operatorname{X}}
\end{aligned}\]

\textbf{Part A: Prove $\operatorname{Z} = \ket{0}\bra{0} - \ket{1}\bra{1} = \ket{+}\bra{-} + \ket{-}\bra{+}$.}

\[\begin{aligned}
\ket{+}\bra{-} &= \left(\frac{1}{\sqrt{2}}(\ket{0} + \ket{1})\right)\left(\frac{1}{\sqrt{2}}(\bra{0} - \bra{1})\right) \\ 
&= \frac{1}{2}(\ket{0}\bra{0} + \ket{0}\bra{1} - \ket{1}\bra{0} - \ket{1}\bra{1})
\end{aligned}\]

\[\begin{aligned}
\ket{-}\bra{+} &= \left(\frac{1}{\sqrt{2}}(\ket{0} - \ket{1})\right)\left(\frac{1}{\sqrt{2}}(\bra{0} + \bra{1})\right) \\ 
&= \frac{1}{2}(\ket{0}\bra{0} - \ket{0}\bra{1} + \ket{1}\bra{0} + \ket{1}\bra{1})
\end{aligned}\]

\[\begin{aligned}
\ket{+}\bra{-} + \ket{-}\bra{+} &= \frac{1}{2}(\ket{0}\bra{0} + \ket{0}\bra{1} - \ket{1}\bra{0} - \ket{1}\bra{1}) + \frac{1}{2}(\ket{0}\bra{0} - \ket{0}\bra{1} + \ket{1}\bra{0} + \ket{1}\bra{1}) \\
&= \ket{0}\bra{0} - \ket{1}\bra{1} \\
&= \boxed{\operatorname{Z}}
\end{aligned}\]

\textbf{Part B: Preparation}

\[\operatorname{I} = \ket{0}\bra{0} + \ket{1}\bra{1}\]

A very useful identity (valid for any single-qubit operator $U$) is the following:

\[\ket{0}\bra{0} \otimes \operatorname{I} + \ket{1}\bra{1} \otimes \operatorname{U}\]

means: “apply $\operatorname{U}$ on the second qubit iff the first qubit is $\ket{1}$”. (control in computational/
$\operatorname{Z}$ basis). Similarly,

\[\ket{+}\bra{+} \otimes \operatorname{I} + \ket{-}\bra{-} \otimes \operatorname{U}\]

means: control in the $\operatorname{X}$-basis. (control is $\ket{-}$).

\textbf{Part B: Prove $\ket{0}\bra{0} \otimes \operatorname{I} + \ket{1}\bra{1} \otimes \operatorname{Z} = \operatorname{I} \otimes \ket{0}\bra{0} + \operatorname{Z} \otimes \ket{1}\ket{1}$.}

\[\begin{aligned}
\text{LHS} &= \ket{0}\bra{0} \otimes \operatorname{I} + \ket{1}\bra{1} \otimes \operatorname{Z} \\
&= \ket{0}\bra{0} \otimes (\ket{0}\bra{0} + \ket{1}\bra{1})+ \ket{1}\bra{1} \otimes (\ket{0}\bra{0} - \ket{1}\bra{1}) \\
&= \ket{0}\bra{0} \otimes \ket{0}\bra{0} + \ket{0}\bra{0} \otimes \ket{1}\bra{1} + \ket{1}\bra{1} \otimes \ket{0}\bra{0} - \ket{1}\bra{1} \otimes \ket{1}\bra{1} \\
&= \boxed{\ket{00}\bra{00} + \ket{01}\bra{01} + \ket{10}\bra{10} - \ket{11}\bra{11}}
\end{aligned}\]

\[\begin{aligned}
\text{RHS} &= \operatorname{I} \otimes \ket{0}\bra{0} + \operatorname{Z} \otimes \ket{1}\bra{1} \\
&= (\ket{0}\bra{0} + \ket{1}\bra{1}) \otimes \ket{0}\bra{0} + (\ket{0}\bra{0} - \ket{1}\bra{1}) \otimes \ket{1}\bra{1} \\
&= \ket{0}\bra{0} \otimes \ket{0}\bra{0} + \ket{1}\bra{1} \otimes \ket{0}\bra{0} + \ket{0}\bra{0} \otimes \ket{1}\bra{1} - \ket{1}\bra{1} \otimes \ket{1}\bra{1} \\
&= \boxed{\ket{00}\bra{00} + \ket{01}\bra{01} + \ket{10}\bra{10} - \ket{11}\bra{11}}
\end{aligned}\]

Therefore, \[\boxed{\text{LHS} = \text{RHS}}\]

\textbf{Part B: Prove $\ket{+}\bra{+} \otimes \operatorname{I} +\ket{-}\bra{-} \otimes \operatorname{X} = \operatorname{I} \otimes \ket{+}\bra{+} + \operatorname{X} \otimes \ket{-}\bra{-}$.}

\[\ket{+}\bra{+} = \frac{1}{2}(\operatorname{I} + \operatorname{X}), \quad \ket{-}\bra{-} = \frac{1}{2}(\operatorname{I} - \operatorname{X})\]

\[\begin{aligned}
\text{LHS} &= \ket{+}\bra{+} \otimes \operatorname{I} +\ket{-}\bra{-} \otimes \operatorname{X} \\
&= \frac{1}{2}(\operatorname{I} + \operatorname{X}) \otimes \operatorname{I} + \frac{1}{2}(\operatorname{I} - \operatorname{X}) \otimes \operatorname{X} \\
&= \frac{1}{2}(\operatorname{I} \otimes \operatorname{I} + \operatorname{X} \otimes \operatorname{I}) + \frac{1}{2}(\operatorname{I} \otimes \operatorname{X} - \operatorname{X} \otimes \operatorname{X}) \\
&= \boxed{\frac{1}{2}(\operatorname{I} \otimes \operatorname{I} + \operatorname{I} \otimes \operatorname{X} + \operatorname{X} \otimes \operatorname{I} - \operatorname{X} \otimes \operatorname{X})}
\end{aligned}\]

\[\begin{aligned}
\text{RHS} &= \operatorname{I} \otimes \ket{+}\bra{+} + \operatorname{X} \otimes \ket{-}\bra{-} \\
&= \operatorname{I} \otimes \frac{1}{2}(\operatorname{I} + \operatorname{X}) + \operatorname{X} \otimes \frac{1}{2}(\operatorname{I} - \operatorname{X}) \\
&= \frac{1}{2}(\operatorname{I} \otimes \operatorname{I} + \operatorname{I} \otimes \operatorname{X}) + \frac{1}{2}(\operatorname{X} \otimes \operatorname{I} - \operatorname{X} \otimes \operatorname{X}) \\
&= \boxed{\frac{1}{2}(\operatorname{I} \otimes \operatorname{I} + \operatorname{I} \otimes \operatorname{X} + \operatorname{X} \otimes \operatorname{I} - \operatorname{X} \otimes \operatorname{X})}
\end{aligned}\]

So \[\boxed{\text{LHS} = \text{RHS}}\]

\textbf{Part B: Quantum Circuit Diagrams}

The figure below shows the quantum circuit equivalences for both equations:

\begin{figure}[H]
\centering
\includegraphics[width=0.95\textwidth]{HW2_code/problem5_all_circuits.png}
\caption{Quantum circuit equivalences for Problem 5(B). \textbf{Top row (Equation 1)}: Controlled-Z gate is symmetric—swapping control and target gives the same operation. \textbf{Bottom row (Equation 2)}: Controlled-X with control in the X-basis (using Hadamard gates) is also symmetric.}
\label{fig:problem5_circuits}
\end{figure}

\end{solution}

\section*{Problem 6 (10 points)}

For $\ket{\psi} = \frac{1}{\sqrt{2}} (1,-1)^T$ and $\ket{\phi} = \frac{1}{\sqrt{3}} (\sqrt{2},i)^T$, find $(\ket{i}\bra{j})\ket{\psi}$ and $(\ket{i}\bra{j})\ket{\phi}$.

\begin{solution}

Given:
\[\ket{\psi} = \frac{1}{\sqrt{2}}\begin{bmatrix}1 \\ -1\end{bmatrix}, \quad \ket{\phi} = \frac{1}{\sqrt{3}}\begin{bmatrix}\sqrt{2} \\ i\end{bmatrix}\]

\[\ket{0} = \begin{bmatrix}1 \\ 0\end{bmatrix}, \quad \ket{1} = \begin{bmatrix}0 \\ 1\end{bmatrix}, \quad \bra{0} = \begin{bmatrix}1 & 0\end{bmatrix}, \quad \bra{1} = \begin{bmatrix}0 & 1\end{bmatrix}\]

\textbf{Part 1: Compute $(\ket{i}\bra{j})\ket{\psi}$ for all $i,j \in \{0,1\}$}

\textbf{Case $j=0$:}

First compute $\bra{0}\ket{\psi}$:
\[\bra{0}\ket{\psi} = \begin{bmatrix}1 & 0\end{bmatrix}\frac{1}{\sqrt{2}}\begin{bmatrix}1 \\ -1\end{bmatrix} = \frac{1}{\sqrt{2}}\]

So:
\[(\ket{0}\bra{0})\ket{\psi} = \ket{0}(\bra{0}\ket{\psi}) = \frac{1}{\sqrt{2}}\ket{0}\]

\[(\ket{1}\bra{0})\ket{\psi} = \ket{1}(\bra{0}\ket{\psi}) = \frac{1}{\sqrt{2}}\ket{1}\]

\textbf{Case $j=1$:}

First compute $\bra{1}\ket{\psi}$:
\[\bra{1}\ket{\psi} = \begin{bmatrix}0 & 1\end{bmatrix}\frac{1}{\sqrt{2}}\begin{bmatrix}1 \\ -1\end{bmatrix} = -\frac{1}{\sqrt{2}}\]

So:
\[(\ket{0}\bra{1})\ket{\psi} = \ket{0}(\bra{1}\ket{\psi}) = -\frac{1}{\sqrt{2}}\ket{0}\]

\[(\ket{1}\bra{1})\ket{\psi} = \ket{1}(\bra{1}\ket{\psi}) = -\frac{1}{\sqrt{2}}\ket{1}\]

\textbf{Summary for $\ket{\psi}$:}

\[\boxed{\begin{aligned}
(\ket{0}\bra{0})\ket{\psi} &= \frac{1}{\sqrt{2}}\ket{0} \\
(\ket{1}\bra{0})\ket{\psi} &= \frac{1}{\sqrt{2}}\ket{1} \\
(\ket{0}\bra{1})\ket{\psi} &= -\frac{1}{\sqrt{2}}\ket{0} \\
(\ket{1}\bra{1})\ket{\psi} &= -\frac{1}{\sqrt{2}}\ket{1}
\end{aligned}}\]

\textbf{Part 2: Compute $(\ket{i}\bra{j})\ket{\phi}$ for all $i,j \in \{0,1\}$}

\textbf{Case $j=0$:}

First compute $\bra{0}\ket{\phi}$:
\[\bra{0}\ket{\phi} = \begin{bmatrix}1 & 0\end{bmatrix}\frac{1}{\sqrt{3}}\begin{bmatrix}\sqrt{2} \\ i\end{bmatrix} = \frac{\sqrt{2}}{\sqrt{3}}\]

So:
\[(\ket{0}\bra{0})\ket{\phi} = \ket{0}(\bra{0}\ket{\phi}) = \frac{\sqrt{2}}{\sqrt{3}}\ket{0}\]

\[(\ket{1}\bra{0})\ket{\phi} = \ket{1}(\bra{0}\ket{\phi}) = \frac{\sqrt{2}}{\sqrt{3}}\ket{1}\]

\textbf{Case $j=1$:}

First compute $\bra{1}\ket{\phi}$:
\[\bra{1}\ket{\phi} = \begin{bmatrix}0 & 1\end{bmatrix}\frac{1}{\sqrt{3}}\begin{bmatrix}\sqrt{2} \\ i\end{bmatrix} = \frac{i}{\sqrt{3}}\]

So:
\[(\ket{0}\bra{1})\ket{\phi} = \ket{0}(\bra{1}\ket{\phi}) = \frac{i}{\sqrt{3}}\ket{0}\]

\[(\ket{1}\bra{1})\ket{\phi} = \ket{1}(\bra{1}\ket{\phi}) = \frac{i}{\sqrt{3}}\ket{1}\]

\textbf{Summary for $\ket{\phi}$:}

\[\boxed{\begin{aligned}
(\ket{0}\bra{0})\ket{\phi} &= \frac{\sqrt{2}}{\sqrt{3}}\ket{0} \\
(\ket{1}\bra{0})\ket{\phi} &= \frac{\sqrt{2}}{\sqrt{3}}\ket{1} \\
(\ket{0}\bra{1})\ket{\phi} &= \frac{i}{\sqrt{3}}\ket{0} \\
(\ket{1}\bra{1})\ket{\phi} &= \frac{i}{\sqrt{3}}\ket{1}
\end{aligned}}\]

\end{solution}

\section*{Problem 7 (25 points)}

Write down the Dyad expressions (i.e., \textit{outer} products) in the computational
basis for the following 1, 2 and 3-qubit gates, respectively. Also write down the
corresponding matrix representations in the computational basis: \\
(1) $\operatorname{H}$, \\
(2) $\operatorname{CNOT}$ where the 2nd qubit is the control, \\
(3) $\operatorname{CZ}$ where the 2nd qubit is the control, \\
(4) SWAP. Note: From its definition, write it dyads in the computational basis, \\
(5) Toffoli gate where the 2nd qubit is the Target.

\textbf{Hint 1}: \textit{Carefully apply the definitions of each gate and write down the operator remembering that each dyad $\ket{\Psi}\bra{\Phi}$ means if $\ket{\Phi}$ is the input, then $\ket{\Psi}$ is the output. For example, Z-gate is defined as: Z flips the sign for $\ket{1}$ state, but leaves $\ket{0}$ alone $\boldsymbol{\rightarrow} \operatorname{Z} = \ket{0}\bra{0} - \ket{1}\bra{1}$. Another example, $\operatorname{H}$ sends $\ket{0}$ to $\ket{+}$ and $\ket{1}$ to $\ket{-} \boldsymbol{\rightarrow} \operatorname{H} = \ket{+}\bra{0} + \ket{-}\bra{1}$. $\operatorname{H}$ sends $\ket{+}\bra{0} + \ket{-}\bra{1} \boldsymbol{\rightarrow} \operatorname{H} = \ket{0}\bra{+} + \ket{1}\bra{-}$. You can
convert either expression to that all in the computational basis.}

\textbf{Hint 2}: \textit{For the multi-qubit gates, first write down as the tensor products of single qubit dyads (or operators), then convert it to the computational basis. For example, Controlled-Z with the first qubit as the Control: $\operatorname{CZ} = \ket{0}\bra{0} \otimes \operatorname{I} + \ket{1}\bra{1} \otimes \operatorname{Z}$. You can change I and Z to the sum of dyads: $\operatorname{I} = \ket{0}\bra{0} + \ket{1}\bra{1}$ and $\operatorname{Z} = \ket{0}\bra{0} - \ket{1}\bra{1}$, and carry out the tensor, which will give $CZ$ as the sum of dyads in the computational basis.}

\begin{solution}

\textbf{(1) Hadamard gate H}

\textbf{Action:} $\operatorname{H}\ket{0} = \ket{+} = \frac{1}{\sqrt{2}}(\ket{0} + \ket{1})$, $\operatorname{H}\ket{1} = \ket{-} = \frac{1}{\sqrt{2}}(\ket{0} - \ket{1})$

\textbf{Dyad expression (using $\ket{\pm}$ basis):}
\[\operatorname{H} = \ket{+}\bra{0} + \ket{-}\bra{1}\]

\textbf{Dyad expression (computational basis):}

Substitute $\ket{+}$ and $\ket{-}$:
\[\begin{aligned}
\operatorname{H} &= \frac{1}{\sqrt{2}}(\ket{0} + \ket{1})\bra{0} + \frac{1}{\sqrt{2}}(\ket{0} - \ket{1})\bra{1} \\
&= \frac{1}{\sqrt{2}}\ket{0}\bra{0} + \frac{1}{\sqrt{2}}\ket{1}\bra{0} + \frac{1}{\sqrt{2}}\ket{0}\bra{1} - \frac{1}{\sqrt{2}}\ket{1}\bra{1}
\end{aligned}\]

\[\boxed{\operatorname{H} = \frac{1}{\sqrt{2}}(\ket{0}\bra{0} + \ket{0}\bra{1} + \ket{1}\bra{0} - \ket{1}\bra{1})}\]

\textbf{Matrix representation:}
\[\boxed{\operatorname{H} = \frac{1}{\sqrt{2}}\begin{bmatrix}1 & 1 \\ 1 & -1\end{bmatrix}}\]

\textbf{(2) CNOT where 2nd qubit is control}

\textbf{Action:} Flip the 1st qubit (target) if the 2nd qubit (control) is $\ket{1}$

\[\ket{00} \to \ket{00}, \quad \ket{01} \to \ket{11}, \quad \ket{10} \to \ket{10}, \quad \ket{11} \to \ket{01}\]

\textbf{Dyad expression (tensor product form):}
\[\operatorname{CNOT} = \operatorname{I} \otimes \ket{0}\bra{0} + \operatorname{X} \otimes \ket{1}\bra{1}\]

\textbf{Dyad expression (computational basis):}

Expand using $\operatorname{I} = \ket{0}\bra{0} + \ket{1}\bra{1}$ and $\operatorname{X} = \ket{0}\bra{1} + \ket{1}\bra{0}$:
\[\begin{aligned}
\operatorname{CNOT} &= (\ket{0}\bra{0} + \ket{1}\bra{1}) \otimes \ket{0}\bra{0} + (\ket{0}\bra{1} + \ket{1}\bra{0}) \otimes \ket{1}\bra{1} \\
&= \ket{0}\bra{0} \otimes \ket{0}\bra{0} + \ket{1}\bra{1} \otimes \ket{0}\bra{0} + \ket{0}\bra{1} \otimes \ket{1}\bra{1} + \ket{1}\bra{0} \otimes \ket{1}\bra{1}
\end{aligned}\]

\[\boxed{\operatorname{CNOT} = \ket{00}\bra{00} + \ket{10}\bra{10} + \ket{01}\bra{11} + \ket{11}\bra{01}}\]

\textbf{Matrix representation (basis order: $\ket{00}, \ket{01}, \ket{10}, \ket{11}$):}
\[\boxed{\operatorname{CNOT} = \begin{bmatrix}1 & 0 & 0 & 0 \\ 0 & 0 & 0 & 1 \\ 0 & 0 & 1 & 0 \\ 0 & 1 & 0 & 0\end{bmatrix}}\]

\textbf{(3) CZ where 2nd qubit is control}

\textbf{Action:} Apply Z to 1st qubit if 2nd qubit is $\ket{1}$. (Note: CZ is symmetric)

\[\ket{00} \to \ket{00}, \quad \ket{01} \to \ket{01}, \quad \ket{10} \to \ket{10}, \quad \ket{11} \to -\ket{11}\]

\textbf{Dyad expression (tensor product form):}
\[\operatorname{CZ} = \operatorname{I} \otimes \ket{0}\bra{0} + \operatorname{Z} \otimes \ket{1}\bra{1}\]

\textbf{Dyad expression (computational basis):}

Expand using $\operatorname{I} = \ket{0}\bra{0} + \ket{1}\bra{1}$ and $\operatorname{Z} = \ket{0}\bra{0} - \ket{1}\bra{1}$:
\[\begin{aligned}
\operatorname{CZ} &= (\ket{0}\bra{0} + \ket{1}\bra{1}) \otimes \ket{0}\bra{0} + (\ket{0}\bra{0} - \ket{1}\bra{1}) \otimes \ket{1}\bra{1} \\
&= \ket{0}\bra{0} \otimes \ket{0}\bra{0} + \ket{1}\bra{1} \otimes \ket{0}\bra{0} + \ket{0}\bra{0} \otimes \ket{1}\bra{1} - \ket{1}\bra{1} \otimes \ket{1}\bra{1}
\end{aligned}\]

\[\boxed{\operatorname{CZ} = \ket{00}\bra{00} + \ket{10}\bra{10} + \ket{01}\bra{01} - \ket{11}\bra{11}}\]

\textbf{Matrix representation:}
\[\boxed{\operatorname{CZ} = \begin{bmatrix}1 & 0 & 0 & 0 \\ 0 & 1 & 0 & 0 \\ 0 & 0 & 1 & 0 \\ 0 & 0 & 0 & -1\end{bmatrix}}\]

\textbf{(4) SWAP gate}

\textbf{Action:} Exchange the two qubits

\[\ket{00} \to \ket{00}, \quad \ket{01} \to \ket{10}, \quad \ket{10} \to \ket{01}, \quad \ket{11} \to \ket{11}\]

\textbf{Dyad expression (computational basis):}
\[\boxed{\operatorname{SWAP} = \ket{00}\bra{00} + \ket{10}\bra{01} + \ket{01}\bra{10} + \ket{11}\bra{11}}\]

\textbf{Matrix representation:}
\[\boxed{\operatorname{SWAP} = \begin{bmatrix}1 & 0 & 0 & 0 \\ 0 & 0 & 1 & 0 \\ 0 & 1 & 0 & 0 \\ 0 & 0 & 0 & 1\end{bmatrix}}\]

\textbf{(5) Toffoli gate where 2nd qubit is Target}

\textbf{Action:} The 1st and 3rd qubits are controls. Flip the 2nd qubit (target) if both controls are $\ket{1}$.

Qubit order: (1st, 2nd, 3rd)

\[\begin{aligned}
\ket{000} &\to \ket{000}, \quad \ket{001} \to \ket{001}, \quad \ket{010} \to \ket{010}, \quad \ket{011} \to \ket{011} \\
\ket{100} &\to \ket{100}, \quad \ket{101} \to \ket{111}, \quad \ket{110} \to \ket{110}, \quad \ket{111} \to \ket{101}
\end{aligned}\]

\textbf{Dyad expression (tensor product form):}
\[\operatorname{Toffoli} = \ket{0}\bra{0} \otimes \operatorname{I} \otimes \operatorname{I} + \ket{1}\bra{1} \otimes \operatorname{I} \otimes \ket{0}\bra{0} + \ket{1}\bra{1} \otimes \operatorname{X} \otimes \ket{1}\bra{1}\]

\textbf{Dyad expression (computational basis):}
\[\begin{aligned}
\operatorname{Toffoli} = &\ket{000}\bra{000} + \ket{001}\bra{001} + \ket{010}\bra{010} + \ket{011}\bra{011} \\
&+ \ket{100}\bra{100} + \ket{101}\bra{111} + \ket{110}\bra{110} + \ket{111}\bra{101}
\end{aligned}\]

\[\boxed{\operatorname{Toffoli} = \ket{000}\bra{000} + \ket{001}\bra{001} + \ket{010}\bra{010} + \ket{011}\bra{011} + \ket{100}\bra{100} + \ket{111}\bra{101} + \ket{110}\bra{110} + \ket{101}\bra{111}}\]

\textbf{Matrix representation (basis order: $\ket{000}, \ket{001}, \ket{010}, \ket{011}, \ket{100}, \ket{101}, \ket{110}, \ket{111}$):}
\[\boxed{\operatorname{Toffoli} = \begin{bmatrix}
1 & 0 & 0 & 0 & 0 & 0 & 0 & 0 \\
0 & 1 & 0 & 0 & 0 & 0 & 0 & 0 \\
0 & 0 & 1 & 0 & 0 & 0 & 0 & 0 \\
0 & 0 & 0 & 1 & 0 & 0 & 0 & 0 \\
0 & 0 & 0 & 0 & 1 & 0 & 0 & 0 \\
0 & 0 & 0 & 0 & 0 & 0 & 0 & 1 \\
0 & 0 & 0 & 0 & 0 & 0 & 1 & 0 \\
0 & 0 & 0 & 0 & 0 & 1 & 0 & 0
\end{bmatrix}}\]

\end{solution}

\end{document}